\section{RTWBS: can a new bisimulation quotient more than \uppaal{}? \hfill \open{}}

RTWBS (Relaxed Weak Timed Bisimulation) is an \emph{asymmetric} refinement relation: it insists
that a refined component preserve the timing of everything it \emph{sends}, while allowing it to be
slower about when it \emph{receives}. It is built on weak timed bisimulation, so internal
($\tau$) steps are abstracted away, and it is decided by partition refinement over the zone graphs
of the two components.

\paragraph{First, the mismatch to be honest about.} RTWBS as published answers ``does this update
refine the deployed system?'' --- a question about \emph{two} systems. Reducing the state space of
\emph{one} system is a different use. To get a quotient you take the kernel of the preorder (relate
$s$ and $t$ when each refines the other) and factor by it. That is legitimate, but it is a
different theorem from the ones in the paper, and it needs stating.

\paragraph{Why it could nevertheless be the best lead in this document.} Three reasons.

\begin{enumerate}
\item \textbf{\uppaal{} does no bisimulation quotienting at all.} It reduces by zone abstraction and
      inclusion. So any working quotient is strictly extra, not a re-implementation of something
      they have.

\item \textbf{Weak is the point.} A \emph{strong} bisimulation quotient on these models would
      collapse almost nothing, for the same reason \S3 fails: every discrete move changes what
      others can do. A \emph{weak} one abstracts internal steps, and here the property names only
      two processes --- so every move of the other $N-2$ is a candidate $\tau$. That is precisely
      the structure the relation is built to exploit, and it attacks the discrete dimension, which
      is the only dimension with slack.

\item \textbf{It composes with \S1.} The environment-abstraction lead needs a small $\hat E$ that is
      at least as permissive as the real environment. A weak-bisimulation quotient of the
      environment is a principled way to \emph{construct} one, instead of writing it by hand ---
      and hand-construction is exactly the step that stalled that line of work in 1997.
\end{enumerate}

\paragraph{The obstacle, stated plainly.} Computing a quotient of the product needs the product,
which is the thing we are trying not to build. The technique only pays if the quotient is taken
\emph{per component, before composition} --- compositional minimisation. Whether the relation is a
congruence for \uppaal{}'s composition (shared bounded integers, binary and broadcast channels,
committed locations) is the question that decides the whole lead, and it is not answered in the
paper, which composes contracts rather than quotients components.

\paragraph{Two smaller things worth noting.} The paper observes that pair validation is
``naturally parallelisable'' with an effective speedup up to the core count, which fits this
engine's setting exactly. And the relation preserves CTL$^*$ over \emph{observable events}, whereas
our properties are predicates over \emph{locations}; those need to be reconciled before any
soundness claim, most likely by making the property's locations observable and everything else
internal.

\paragraph{What would settle it, cheaply.} Take \texttt{fischer}, designate every action of
processes $3..N$ as internal, compute the weak-bisimulation quotient of a \emph{single} process
template, and count the resulting classes. If one template collapses, compositional minimisation is
worth building. If it does not, the lead dies for the same reason \S3 died, and we will know in an
afternoon.
